\section{A concrete cross-protocol execution trace}\label{sec:C19:trace}
Artifact transfer needs an explicit trust boundary. An A2A artifact identifier belongs to the producing task's namespace; an MCP resource identifier belongs to the serving endpoint's namespace. Resolving one into the other is an application mapping, not an equality implied by the identifier strings. A content digest can detect a changed payload but cannot establish the producer's authority or the truth of its claims. Preserve the producing task, source version, content type, and verification status when constructing the resource representation. The consumer should distinguish an agent's inference from evidence that the remote agent actually observed. Otherwise a technically correct transfer can launder a speculative conclusion into an apparently authoritative tool result.

An interrupted task also needs a scope-preserving continuation rule. Suppose a remote agent requests authorization for one operation and later resumes with changed arguments. Reusing the earlier approval solely because the task identifier is unchanged can authorize a different effect. A proposed application contract binds a grant to the relevant argument digest, resource scope, authority epoch, and expiry; a changed material precondition requires revalidation. This is an authority requirement derived from the execution model, not a claim that every A2A implementation enforces that binding. Protocol progress and user authorization are separate state machines, even when one user interface displays them together.

Consider a constructed workflow in which a coordinator delegates a research task to a remote agent and later uses a tool server to retrieve a source record. The A2A request identifies the objective, current constraints, and expected evidence artifact. The remote service returns a task handle or a direct message according to its supported interaction. The coordinator records the handle in durable task state rather than relying on an open stream to remember it.

When evidence requires a tool server, the coordinator's MCP client sends a request carrying the current protocol version and client capabilities. It may first use discovery, then list the authorized tools and obtain the selected schema. Before invoking the tool, the host validates arguments and task authority. The server's response becomes a typed observation with server and resource identity. It does not become a new user instruction merely because it was returned through a trusted transport.

Suppose the research task returns a completed artifact while one cited resource has changed. The coordinator should validate the cited version or preserve the captured passage, then decide whether the artifact satisfies the requested evidence standard. If the artifact proposes publication but the delegation allowed only preparation, local authority remains binding. A2A task completion and MCP tool success do not jointly imply permission to publish.

Now interrupt the network after a tool operation may have committed but before the coordinator receives the result. Reconnecting with the correct protocol metadata does not resolve the effect. The application needs an explicit operation identifier, remote status query, or reconciliation procedure. If no such mechanism exists, it should retain an outcome-unknown state. Reissuing the same semantic request with a new JSON-RPC identifier can duplicate work unless the tool supplies an independent idempotency contract.

This trace is original synthesis, not a tested integration or a claim that every SDK implements these steps. It identifies the records a correct implementation would need: delegation identifier, task/context relation, protocol versions, catalog/schema digest, authorization version, effect identifier, response receipt, and artifact acceptance decision. These records support causal diagnosis if the final output is wrong.

The most informative interoperability experiment would replay this workflow with independent client and server implementations, controlled delays, credential changes, and a mixed modern/legacy MCP environment. A successful result would support the tested semantic contract, not merely that two programs exchanged JSON. A failed result should identify whether the mismatch concerned version negotiation, capability declaration, effect semantics, or the meaning of task completion.

The normative implementation anchor is the versioned TypeScript schema. At the inspected July 28 commit, \texttt{RequestMetaObject} requires the protocol-version and client-capabilities members while declaring client information optional. \texttt{DiscoverRequest} fixes the discovery method, and \texttt{CallToolRequest} binds the tool-call method to its parameter structure~\cite{mcp2026schema}. These definitions are stronger implementation references than copying an old SDK example and assuming it speaks the current protocol. They describe message shape; they are not an executed conformance test of any client library.

A concrete negative test removes only the capabilities member from the illustrative request while preserving valid JSON and a recognized version. A second test keeps the capabilities member but changes the HTTP version header so it disagrees with the body. A third requests an unsupported protocol version while keeping body and header equal. These cases should not collapse into one generic retry: malformed metadata, inconsistent transport declaration, and unsupported version are different failures. Their distinction matters because only some errors justify selecting a different protocol version.

An application can generate clients from schema definitions and still mishandle those outcomes. It might cache a successful discovery response under one credential and reuse the catalog under another. It might regard an optional identity string as a trusted principal. It might strip required metadata when forwarding a request to a worker. These are integration failures above the type definition. Test the complete forwarding path with different principals and concurrent tasks rather than validating only an isolated client serializer.

For long-running operations, retain two kinds of correlation. The protocol request identifier tracks a particular exchange, while an explicit application or extension handle tracks durable work. After a reconnect, a new exchange can refer to the same durable task. If the implementation uses request identity as the only operation identity, it can either duplicate work on retry or incorrectly suppress a legitimate new request. A reference fixture should model both identities and assert which one is stable across network retries.

The same distinction supports audit of delegated research. The coordinator can record that a remote task completed and that an artifact was received, while separately recording whether the artifact passed evidence validation. This avoids turning transport success into scientific acceptance. It also allows the coordinator to reject one artifact without falsely marking the remote operation as never completed. Protocol state, task state, and epistemic status remain separate even when the user interface displays one compact progress indicator.

